\section{Main results}
\label{sec:mainresult}

Let $\Pi$ be a probability distribution representing our prior belief about the effect size. Assume it admits a density $\pi_m$ with respect to the Lebesgue measure, where $m\in\bbR$ is the known center of the distribution. Assume it admits the location-shift construction $\pi_m(\theta)=\pi(\theta-m)$, for some density $\pi$.\fc{should add more discussion to justify the limit taken} 

For \(u\in(-\sqrt n,\sqrt n)\), define
\[
h_n(\theta,u)
=
e^{-n\theta^2/2}A_n(u\theta),
\]
and
\[
M_{n,m}(u)
=
\int_{\mathbb{R}}
h_n(\theta,u)\pi(\theta-m)\,d\theta,
\qquad
M_n^{\Leb}(u)
=
\int_{\mathbb{R}}h_n(\theta,u)\,d\theta.
\]
Thus \(M_{n,m}=M_{n,m}(U_n)\) and
\(M_n^{\Leb}=M_n^{\Leb}(U_n)\). Also define
\[
g_{\Leb}(c)
=
\mathbb{E}_0\left[
\min\left(\frac{M_2^{\Leb}(U_2)}{c},1\right)
\right],
\qquad c>0.
\]


%
%Define
%\begin{equation}\label{eq:Mnm-def}
%  M_{n,m}\;=\;\int_\mathbb{R} e^{-n\theta^2/2}A_n(\theta U_n)\,\pi_m(\theta)
%  \,\dd\theta,
%  \qquad
%  g_m(c)\;=\;\E_0\!\bigl[\min(M_{2,m}/c,\,1)\bigr].
%\end{equation}
We are interested in the behaviour of the test martingale and associated $p$-processes, for a given sample size $n$, when the prior becomes badly misspecified. That is, for any fixed $n$ and any fixed dataset $y_{1:n}$, when the centering $|m|\to \infty$. We first prove a negative result: that Ville-calibrated $p$-variables become powerless in case of a bad prior. 


\begin{proposition}
For any proper prior $\Pi$, any $n\geq 2$ and any fixed dataset $y_{1:n}$ with \(|u_n|=\frac{|\sum_{i=1}^ny_i|}
       {\sqrt{\sum_{i=1}^n y_i^2}}<\sqrt n\), we have, as $m\to\pm\infty$,
 $$M_{n,m}\to 0,$$
 hence the Ville-calibrated  $p$-process satisfies, as $m\to \pm \infty$,
$$
P^{V}_{n,m}:=\min(1,M_{n,m}^{-1})\to 1.
$$
\end{proposition}
The above result holds for any dataset; that is, however strong the evidence against null is, the Ville-calibrated $p$-variable becomes powerless when the prior becomes badly centered. Note that the result excludes the extreme cases $|u_n|=\sqrt{n}$, where all the observations take the same value. We note that the same result holds for $\widetilde P_{n,m}^V$.


\begin{proof}
Let \(Z\) have density \(\pi\). By the location-shift
representation \(\pi_m(\theta)=\pi(\theta-m)\), the change of
variables \(\theta=m+z\) gives
\[
M_{n,m}
=
\mathbb{E}\left[h_n(m+Z,u_n)\right],
\]
where
\[
h_n(\theta,u)
=
e^{-n\theta^2/2}A_n(u\theta).
\]

Since \(|u_n|<\sqrt n\), choose \(\eta\) such that
\[
u_n^2<\eta<n.
\]
By Lemma~\ref{lemma:hthetabound}, there exist constants
\(C_{n,u_n},\kappa_{n,u_n}>0\) such that
\[
0\leq h_n(\theta,u_n)
\leq
C_{n,u_n}e^{-\kappa_{n,u_n}\theta^2},
\qquad \theta\in\mathbb{R}.
\]
Consequently, for every fixed \(z\in\mathbb{R}\),
\[
h_n(m+z,u_n)\longrightarrow0
\qquad\text{as }m\to\pm\infty,
\]
and
\[
0\leq h_n(m+z,u_n)\leq C_{n,u_n}.
\]
Since the dominating constant is integrable with respect to the
probability distribution of \(Z\), the dominated convergence
theorem yields
\[
M_{n,m}
=
\mathbb{E}\left[h_n(m+Z,u_n)\right]
\longrightarrow0.
\]
Because \(M_{n,m}>0\), it follows that
\[
P_{n,m}^{V}
=
\min\left(1,M_{n,m}^{-1}\right)
\longrightarrow1.
\]
\end{proof}

%\begin{proof}
%For $Z\sim \pi$, we have
%$$
%M_{n,m}=\bbE[h_{n}(m+Z,u_n)]
%$$
%where
%$$
%h_{n}(\theta,u)=e^{-n\theta^2/2 }A_n(u \theta ).
%$$
%
%%
%%where $X\sim\chi_n$. For every $x\in\bbR$ and $X\geq 0$,
%%$$
%%-nx^2/2 +uxX\leq \frac{u^2 X^2}{2n}.
%%$$
%%Therefore, for $|u|<\sqrt{n}$ and every $x$, we have the uniform (in $x$) bound
%%$$
%%0\leq h_u(x)\leq \bbE[e^{u^2 X^2/(2n)}]=\left(1-\frac{u^2}{n}\right)^{-n/2}.
%%$$
%For $t>0$, the MGF of $\chi_n$ satisfies
%$$
%A_n(t)\sim \frac{\sqrt{2\pi}}{2^{n/2-1}\Gamma(n/2)}t^{n-1}e^{t^2/2}.
%$$
%Consequently, for $|u_n|<\sqrt{n}$, and any $z$, $h_{n,m+z}(u_n)\to 0$ as $m\to\pm\infty$. Additionally, by Lemma \ref{lemma:hthetabound}, $\theta\mapsto h_{n}(\theta,u)$ is bounded for any $n\geq 2$ and $|u|<\sqrt{n}$. Therefore, by dominated convergence, $M_{n,m}\to 0$.
%\end{proof}





\begin{assumption}\label{ass:tails}
Assume that $\pi$ has regularly varying tails: for some $\beta>1$, some slowly varying functions $L_{+}$ and $L_-$,
\begin{align}
\pi(x)\sim |x|^{-\beta}L_{\pm} (|x|),\qquad x\to\pm\infty.
\end{align}
\end{assumption}

Examples include the Student-t with $\nu$ degrees of freedom ($\beta=\nu+1$) the Cauchy ($\beta=2$), and the horseshoe location family ($\beta=2$). In all these example, the slowly varying functions are positive constants.

Set
\(a_m =\pi_m(0)=\pi(-m)\sim m^{-\beta}L_-(m)\).


\begin{theorem}\label{thm:recovery}
Fix \(n\geq 2\) and a realization \(y_{1:n}\) with
\(|u_n|=\frac{|\sum_{i=1}^ny_i|}
       {\sqrt{\sum_{i=1}^n y_i^2}}<\sqrt n\). Under
Assumption~\ref{ass:tails}, as \(m\to\infty\),
\begin{align}
  \frac{M_{n,m}}{a_m}
   &\;\longrightarrow\; M_n^{\Leb},
   \label{eq:Mnm-asymp}\\
  g_m(a_m c)
   &\;\longrightarrow\; g_{\Leb}(c),\qquad c>0,
   \label{eq:gm-asymp}\\
  P_{n,m}^E=g_m\!\bigl(M_{n,m}\bigr)
   &\;\longrightarrow\;
   g_{\Leb}\!\bigl(M_n^{\Leb}\bigr)
   \;=\;P_n^{\Leb}.
   \label{eq:p-asymp}
\end{align}
\end{theorem}

%
%\fc{to be moved later}
%
%\begin{proof}
%We recall that
%$$
%h_{n}(\theta,u)=e^{-n\theta^2/2 }A_n(u \theta ).
%$$
%Write $M_{n,m}(u)=\int h_n(\theta,u)\Pi(\dd\theta)$ and $M_{n,m}^\Leb(u)=\int h_n(\theta,u)\dd\theta$.
%
%\emph{Step 1 (ratio convergence).}  Fix $R>0$ and split
%\[
%  \frac{M_{n,m}(u)}{a_m}
%  \;=\;\underbrace{\int_{|\theta|\leq R}}_{=:I_1(R,m)}
%   +\underbrace{\int_{R<|\theta|\leq m/2}}_{=:I_2(R,m)}
%   +\underbrace{\int_{|\theta|>m/2}}_{=:I_3(m)}\;
%   h_{n}(\theta,u)\,\frac{\pi(\theta-m)}{a_m}\,\dd\theta.
%\]
%We have, for any $\theta$,
%$$
%\pi(\theta-m)\sim a_m \sim m^{-\beta}L_-(m)
%$$
%as $m\to\infty$. By the uniform convergence theorem for regularly varying functions
%\citep[Theorem~1.2.1]{Bingham1987},
%$$
%\sup_{|\theta|\leq R}|\pi(\theta-m)/a_m-1|\to 0
%$$
%as $m\to\infty$, so
%$$
%I_1(R,m)\to\int_{|\theta|\leq R}h_{n}(\theta,u)\,\dd\theta.
%$$
%For $I_2$, on $|\theta|\leq m/2$ we have $|\theta-m|\in[m/2,\,3m/2]$, so
%regular variation yields a constant $C_\star<\infty$ (independent of $R$ and of
%$m$ for all large $m$) such that $\pi(\theta-m)/a_m\leq C_\star$; hence
%$$
%I_2(R,m)\leq C_\star\int_{|\theta|>R}h_{n}(\theta,u)\,\dd\theta,
%$$ which vanishes
%as $R\to\infty$. For $I_3$, using the fact that $\pi$ is a proper prior
%(i.e.\ $\int_\mathbb{R}\pi=1$),
%\[
%I_3(m)\;=\;\int_{|\theta|>m/2}h_{n}(\theta,u)\,\frac{\pi(\theta-m)}{a_m}\,\dd\theta
%   \;\leq\; a_m^{-1}\sup_{|\theta|>m/2}h_{n}(\theta,u),
%\]
%the supremum is exponentially small in $m$ by \eqref{eq:hu-env}, while $a_m$ is
%only polynomially small, so $I_3(m)\to 0$.
%
%Combining these three bounds, for each fixed $R>0$,
%\begin{align*}
%   \int_{|\theta|\leq R}h_{n}(\theta,u)\,\dd\theta
%   \;\leq\; \liminf_{m\to\infty}\frac{M_{n,m}(u)}{a_m}
%   &\;\leq\; \limsup_{m\to\infty}\frac{M_{n,m}(u)}{a_m}\\
%   & \;\leq\; \int_{|\theta|\leq R}h_{n}(\theta,u)\,\dd\theta
%   \;+\; C_\star\!\int_{|\theta|>R}h_{n}(\theta,u)\,\dd\theta.
%\end{align*}
%Sending $R\to\infty$ and using $\int_\mathbb{R} h_{n}(\theta,u)\,\dd\theta
%= M_n^{\Leb}(u)$ gives \eqref{eq:Mnm-asymp}.
%
%\emph{Step 2 (calibration convergence).}
%Since \(|U_2|<\sqrt 2\) almost surely under \(\Prob_0\), Step~1 applied with \(n=2\) gives
%\(M_{2,m}(U_2)/a_m\to M_2^{\Leb}(U_2)\) almost surely. Since \(0\leq\min(M_{2,m}(U_2)/(a_m c),1)\leq 1\), bounded convergence yields
%\[
%  g_m(a_m c)
%  \;=\;\E_0\!\Bigl[\min\!\Bigl(\tfrac{M_{2,m}(U_2)}{a_mc},1\Bigr)\Bigr]
%  \;\longrightarrow\;
%  \E_0\!\Bigl[\min\!\Bigl(\tfrac{M_2^{\Leb}(U_2)}{c},1\Bigr)\Bigr]
%  \;=\;g_{\Leb}(c),
%\]
%proving \eqref{eq:gm-asymp}.
%
%\emph{Step 3 (composition).}
%Set \(c_m:=M_{n,m}(u_n)/a_m\) and \(c:=M_n^{\Leb}(u_n)\); by Step~1 \(c_m\to c\). Each \(g_m\) is decreasing, so for any \(\varepsilon\in(0,c)\) and \(m\) large, \(g_m(a_m(c+\varepsilon))\leq g_m(a_m c_m)\leq g_m(a_m(c-\varepsilon))\). Taking \(\liminf\) and \(\limsup\) using Step~2 and letting \(\varepsilon\downarrow 0\) yields \eqref{eq:p-asymp}.
%\end{proof}




That is, for every fixed \(n\geq 2\) and any dataset with positive sample variance, the extended-Ville $p$-variable under the shifted polynomial-tail prior converges to the extended-Ville $p$-variable under the noninformative Lebesgue prior.


